\documentclass[10pt]{article}
\usepackage[T1]{fontenc}
\usepackage[utf8]{inputenc}
\usepackage[french]{babel}
\usepackage{amsmath,amssymb,mathrsfs,mathpazo}
\usepackage{geometry}
\geometry{a4paper,margin=2cm}
\usepackage{graphicx,array,tabularx,tabularray,colortbl,multirow,multicol,booktabs,enumitem,pifont,ulem,fancybox,diagbox,cancel}
\usepackage{pstricks,pstricks-add,pst-plot,pst-tree,pst-node,pst-text,pst-func,pst-math,pst-bspline,pst-3dplot,pst-eucl,pst-solides3d}
\let\pscalculate\undefined
\usepackage{pst-bezier}
\usepackage{tikz,pgfplots,tkz-tab}
\usetikzlibrary{arrows,arrows.meta,patterns,calc,positioning,trees,decorations.markings,shapes.geometric,shapes.misc,shapes.arrows,intersections,decorations.pathmorphing,decorations.pathreplacing,backgrounds,shadows,babel}
\pgfplotsset{compat=1.18}
\usepgfplotslibrary{fillbetween}
\usepackage[np]{numprint}
\usepackage{listings,makecell,pas-tableur,twoopt}
\let\barre\undefined
\usepackage[pstricks]{tablvar}
\definecolor{darkgreen}{RGB}{0,100,0}
\providecommand{\red}{\color{red}}
\providecommand{\blue}{\color{blue}}
\providecommand{\green}{\color{green}}
\providecommand{\black}{\color{black}}
\providecommand{\white}{\color{white}}
\providecommand{\gray}{\color{gray}}
\providecommand{\cyan}{\color{cyan}}
\providecommand{\magenta}{\color{magenta}}
\providecommand{\yellow}{\color{yellow}}
\providecommand{\violet}{\color{violet}}
\providecommand{\orange}{\color{orange}}
\providecommand{\euro}{\texteuro}
\providecommand{\EUR}{\texteuro}
\providecommand{\ieme}{\textsuperscript{e}}
\providecommand{\degres}{\ensuremath{^\circ}}
\providecommand{\degree}{\ensuremath{^\circ}}
\providecommand{\vv}[1]{\overrightarrow{#1}}
\providecommand{\Coord}[1]{#1}
\providecommand{\decosix}{}
\providecommand{\starredbullet}{\ensuremath{\bullet}}
\pagestyle{empty}
\setlength{\parindent}{0pt}
\usepackage{adjustbox}\DeclareUnicodeCharacter{2236}{:}\DeclareUnicodeCharacter{2190}{\ensuremath{\leftarrow}}
\providecommand{\R}{\mathbb{R}}
\providecommand{\N}{\mathbb{N}}
\providecommand{\Z}{\mathbb{Z}}
\providecommand{\Q}{\mathbb{Q}}
\providecommand{\D}{\mathbb{D}}
\providecommand{\C}{\mathbb{C}}
\providecommand{\e}{{\mathrm{e}}}
\providecommand{\vect}[1]{\overrightarrow{#1}}
\providecommand{\vectt}[1]{\overrightarrow{\mathrm{#1}}}
\providecommand{\vvt}[1]{\overrightarrow{\mathrm{#1}}}
\providecommand{\np}[2][]{{#2~#1}}
\providecommand{\vv}[1]{{\overrightarrow{#1}}}
\renewcommand{\d}{{\mathrm{d}}}
\providecommand{\degres}{{^\circ}}
\providecommand{\degre}{{^\circ}}
\providecommand{\blue}{\color{blue}}
\providecommand{\red}{\color{red}}
\providecommand{\enskip}{\quad}
\providecommand{\llbracket}{\lbrack\!\lbrack}
\providecommand{\rrbracket}{\rbrack\!\rbrack}
\providecommand{\euro}{{\text{€}}}
\providecommand{\no}{{\text{n°}}}
\providecommand{\interval}[2]{{[#1;#2]}}
\providecommand{\intervFO}[2]{{[#1;#2[}}
\providecommand{\dotfill}{\ldots}

\providecommand{\ds}{\displaystyle}
\providecommand{\textding}[1]{#1}
\providecommand{\Oij}{\(\mathrm{O};\vec{\imath},\vec{\jmath}\)}
\providecommand{\Oijk}{\(\mathrm{O};\vec{\imath},\vec{\jmath},\vec{k}\)}
\providecommand{\Ouv}{\(\mathrm{O};\vec{u},\vec{v}\)}
\providecommand{\intervFF}[2]{\ensuremath{[#1;#2]}}
\renewcommand{\euro}{\ifmmode\text{\texteuro}\else\texteuro\fi}
\DeclareUnicodeCharacter{00B0}{\ifmmode^{\circ}\else\textdegree\fi}
\DeclareUnicodeCharacter{20AC}{\ifmmode\text{\texteuro}\else\texteuro\fi}
% Énoncés : Éducation nationale. Transcription et corrigés : APMEP, Denis Vergès et contributeurs. Découpage et classement : CoopMaths, Sébastien Lozano, Éric Elter et contributeurs. Adaptation de la présentation : Maths Explorer.
% CC BY-SA : https://coopmaths.fr/www/about/
% https://coopmaths.fr/alea/static/bac/2023/tex/bac_2023_03_sujet2_reunion_1_cor.tex
\begin{document}
\noindent\textit{APMEP / CoopMaths. CC BY-SA. Présentation adaptée par Maths Explorer.}\par\medskip
\medskip

Un commerçant vend deux types de matelas: matelas RESSORTS et matelas MOUSSE. 

On suppose que chaque client achète un seul matelas.

On dispose des informations suivantes :
\begin{itemize}
\item[$\bullet~~$] 20\,\% des clients achètent un matelas RESSORTS. 

Parmi eux, 90\,\% sont satisfaits de leur achat.
\item[$\bullet~~$] 82\,\% des clients sont satisfaits de leur achat.
\end{itemize}



\medskip

\textbf{Partie A}

\medskip








On note $x = P_{\overline{R}}(S)$, où $P_{\overline{R}}(S)$ désigne la probabilité de $S$ sachant que $R$ n'est pas réalisé.

\begin{minipage}{0.58\linewidth}
\begin{enumerate}
\item Recopier et compléter l'arbre pondéré ci-contre décrivant la situation.
\end{enumerate}

\end{minipage}\hfill
\begin{minipage}{0.38\linewidth}
\begin{center}
\pstree[treemode=R,nodesepA=0pt,nodesepB=2.5pt,treesep = 1cm,levelsep=2.5cm]{\TR{}}
{\pstree{\TR{$R$~~}\taput{0,20}}
{\TR{$S$}\taput{0,9}
\TR{$\overline{S}$}\tbput{0,1}
}
\pstree{\TR{$\overline{R}$~~}\tbput{0,80}}
{\TR{$S$}\taput{$x$}
\TR{$\overline{S}$}\tbput{$1 - x$}
}}
\end{center}
\end{minipage}

\begin{enumerate}[start=2]
\item 
D'après la loi des probabilités totales :
$P(S) = P(R \cap S) + P\left(\overline{R} \cap S \right)$ Or :

$\bullet~~P(R \cap S) = P(R) \times P_R(S) = 0,2 \times 0,9 = 0,18$ et 

$\bullet~~P\left(\overline{R} \cap S \right) = P\left(\overline{R}\right) \times P_{\overline{R}}(S) = 0,8 \times x = 0,8x$ et comme $P(S) = 0,82$, on a donc l'équation :

$0,18 + 0,8x = 0,82 \iff 0,8x = 0,64 \iff x = 0,8$.
\item On choisit un client satisfait de son achat.

Quelle est la probabilité qu'il ait acheté un matelas RESSORTS ?


Il faut calculer $P_S(R) = \dfrac{P(S \cap R)}{P(S)} = \dfrac{0,18}{0,82} = \dfrac{18}{82} = \dfrac{9}{41} \approx 0,219$, soit 0,22 au centième près.
\end{enumerate}
\medskip

\medskip

\textbf{Partie B}

\medskip

\begin{enumerate}
\item On choisit 5 clients au hasard.

On considère la variable aléatoire $X$ qui donne le
nombre de clients satisfaits de leur achat parmi ces 5 clients.
	\begin{enumerate}
		\item 
Les paramètres de la loi binomiale sont $ n = 5 $ \quad et \quad  $ p=\np{0.82}$.
		\item 
		

Il faut trouver $P(X \leqslant 3) \approx \np{0,2223}$, d'après la calculatrice soit 0,222 au millième près.
	\end{enumerate}
\item 


	\begin{enumerate}
		\item On note $p_n$ la probabilité que les $n$ clients soient tous satisfaits de leur achat.


La probabilité qu'un acheteur soit satisfait est égale à 0,82, la probabilité que deux acheteurs soient satisfaits est $0,82^2$, \ldots, la probabilité que les $n$ acheteurs soient satisfaits est égale à $0,82^n$.
		\item 
		

Il faut donc résoudre dans $\N$ l'inéquation :

$p_n < 0,01 \iff 0,82^n < 0,01 \iff n \ln 0,82 < \ln 0,01$ par croissance de la fonction logarithme népérien soit encore 

$n > \dfrac{\ln 0,01}{\ln 0,82}$. (car $\ln 0,82 < 0$). Or $\dfrac{\ln 0,01}{\ln 0,82} \approx 23,2$.

Il faut donc au minimum 24 acheteurs.

Moins de 1\,\% des acheteurs seront tous satisfaits si leur nombre dépasse 23.
	\end{enumerate}
\end{enumerate}

\bigskip

\end{document}
